\section{ゲージ理論の一般構造}

\noindent\textbf{この章の中心命題}
\[
\boxed{
D_\mu=\partial_\mu+igA_\mu^aT^a,
\qquad
[D_\mu,D_\nu]=igF_{\mu\nu}^aT^a
}
\]
ゲージ理論では、時空点ごとに内部自由度の基底を取り替えても物理が変わらないことを要請し、その要請からゲージ場、共変微分、場の強さが定まる。

QED では、電子場の位相を時空点ごとに変える局所 $U(1)$ 変換を許すために、電磁場 $A_\mu$ と共変微分 $D_\mu=\partial_\mu+ieA_\mu$ を導入した。この章では、その構造を一般の内部対称性へ拡張する。ただしここで扱うのは、標準模型に進むために必要な最小限の道具立てであり、群論そのものを独立に展開することはしない。

\subsection{内部自由度と局所変換}
場 $\psi(x)$ が複数の内部成分をもつとする。たとえば
\[
\psi(x)=
\begin{pmatrix}
\psi_1(x)\\
\psi_2(x)
\end{pmatrix}
\]
のような二成分場なら、各時空点で二つの成分を混ぜる変換を考えられる。変換を行列 $U$ で表すと
\[
\psi(x)\mapsto U\psi(x)
\]
である。$U$ が全ての時空点で同じなら大域変換であり、$U=U(x)$ のように時空点ごとに変えられるなら局所変換である。

局所変換では、通常の微分は単純には変換しない。実際、
\[
\partial_\mu\psi(x)
\mapsto
\partial_\mu\{U(x)\psi(x)\}
=
U(x)\partial_\mu\psi(x)+(\partial_\mu U(x))\psi(x)
\]
となり、最後の項が余分に現れる。したがって、局所変換の下で理論の形を保つには、通常の微分 $\partial_\mu$ を別の微分に置き換える必要がある。

\subsection{共変微分とゲージ場}
内部対称性の生成子を $T^a$ とし、その交換関係を
\[
[T^a,T^b]=if^{abc}T^c
\]
と書く。$f^{abc}$ は構造定数であり、生成子同士がどのように交換しないかを表している。$U(1)$ の場合には生成子が一つしかないので、この交換関係は実質的に消える。

局所変換に対して同じ形で変換する微分を作るため、各生成子に対応するゲージ場 $A_\mu^a(x)$ を導入し、
\[
D_\mu:=\partial_\mu+igA_\mu^aT^a
\]
と定義する。ここで $g$ は相互作用の強さを表す結合定数である。ゲージ場は、局所変換で生じる余分な項を打ち消すように変換する量である。このとき
\[
\psi(x)\mapsto U(x)\psi(x)
\qquad\Longrightarrow\qquad
D_\mu\psi(x)\mapsto U(x)D_\mu\psi(x)
\]
となるように $A_\mu^a$ の変換則を定める。

したがって、自由場のラグランジアンに現れる $\partial_\mu\psi$ を $D_\mu\psi$ に置き換えれば、局所変換の下で形が保たれる。これは「相互作用項を思いつきで足す」のではなく、「局所的な内部基底の取り替えを物理的に同等な記述とみなす」という原理から相互作用を作る手順である。

\subsection{場の強さ}
ゲージ場そのものの運動を記述するには、電磁場での
\[
F_{\mu\nu}=\partial_\mu A_\nu-\partial_\nu A_\mu
\]
に対応する量が必要である。一般のゲージ理論では、これは共変微分の交換子から定義される。
\[
[D_\mu,D_\nu]=igF_{\mu\nu}^aT^a.
\]
実際に $D_\mu=\partial_\mu+igA_\mu^aT^a$ を代入して計算すると、
\[
F_{\mu\nu}^a
=
\partial_\mu A_\nu^a-\partial_\nu A_\mu^a
-gf^{abc}A_\mu^bA_\nu^c
\]
を得る。最後の項は、生成子が互いに交換しないことから生じる。$U(1)$ のように $f^{abc}=0$ ならこの項は消え、通常の電磁場の形に戻る。

この違いは物理的にも重要である。$U(1)$ ゲージ理論である QED では、光子は電荷をもたないため、光子同士は直接には結合しない。一方、非可換ゲージ理論では $f^{abc}\neq0$ であるため、ゲージ場自身がその相互作用の荷を担う。したがってゲージ粒子同士の自己相互作用が現れる。

\subsection{ゲージ理論のラグランジアン}
フェルミオン場 $\psi$ とゲージ場 $A_\mu^a$ から作る基本的なラグランジアンは
\[
\mathcal{L}
=
\bar{\psi}(i\gamma^\mu D_\mu-m)\psi
-\frac{1}{4}F_{\mu\nu}^aF^{a\mu\nu}
\]
である。第一項は物質場の運動とゲージ場との結合を表し、第二項はゲージ場自身の運動を表す。

QED はこの式でゲージ群を $U(1)$ とした場合である。強い相互作用を記述する QCD は、ゲージ群を $SU(3)$ とし、三種類の色自由度をもつクォーク場にこの構造を適用した理論である。弱い相互作用と電磁相互作用を合わせた電弱理論では、$SU(2)\times U(1)$ のゲージ構造を用いる。

このように、ゲージ理論は「相互作用を媒介する粒子を追加する」という経験的な記述ではなく、「局所的な内部変換を許しても物理が変わらない」という要請を数式化したものである。その要請が、共変微分、ゲージ場、相互作用項の形を同時に決める。
